% Erstellt von Maximilian Nöthe, <maximilian.noethe@tu-dortmund.de>
% ausgelegt für lualatex und Biblatex mit biber

% Kompilieren mit 
% lualatex dateiname.tex
% biber dateiname.bcf
% lualatex dateiname.tex
% lualatex dateiname.tex
% oder einfach mit:
% make

\documentclass[
  tucolor,
  BCOR=12mm,     % 12mm binding corrections, adjust to fit your binding
  parskip=half,  % new paragraphs start with half line vertical space
  open=any,      % chapters start on both odd and even pages
  cleardoublepage=plain,  % no header/footer on blank pages
]{tudothesis}


% Warning, if another latex run is needed
\usepackage[aux]{rerunfilecheck}

% just list chapters and sections in the toc, not subsections or smaller
\setcounter{tocdepth}{1}

%------------------------------------------------------------------------------
%------------------------------ Sprache und Schrift: --------------------------
%------------------------------------------------------------------------------
\usepackage{fontspec}
\defaultfontfeatures{Ligatures=TeX}  % -- becomes en-dash etc.

% german language
\usepackage{polyglossia}
\setdefaultlanguage{english}

% for english abstract and english titles in the toc
%\setotherlanguages{english}

% intelligent quotation marks, language and nesting sensitive
\usepackage[autostyle]{csquotes}

% microtypographical features, makes the text look nicer on the small scale
\usepackage{microtype}

%------------------------------------------------------------------------------
%------------------------ Für die Matheumgebung--------------------------------
%------------------------------------------------------------------------------

\usepackage{amsmath}
\usepackage{amssymb}
\usepackage{mathtools}
\usepackage{slashed}

% Enable Unicode-Math and follow the ISO-Standards for typesetting math
\usepackage[
  math-style=ISO,
  bold-style=ISO,
  sans-style=italic,
  nabla=upright,
  partial=upright,
]{unicode-math}
\setmathfont{Latin Modern Math}

% nice, small fracs for the text with \sfrac{}{}
\usepackage{xfrac}  


%------------------------------------------------------------------------------
%---------------------------- Numbers and Units -------------------------------
%------------------------------------------------------------------------------

\usepackage[
  locale=DE,
  separate-uncertainty=true,
  per-mode=symbol-or-fraction,
]{siunitx}
\sisetup{math-micro=\text{µ},text-micro=µ}

%------------------------------------------------------------------------------
%-------------------------------- tables  -------------------------------------
%------------------------------------------------------------------------------

\usepackage{booktabs}       % stellt \toprule, \midrule, \bottomrule

%------------------------------------------------------------------------------
%-------------------------------- graphics -------------------------------------
%------------------------------------------------------------------------------

\usepackage{graphicx}
\usepackage{grffile}
\usepackage{tikz}
\usetikzlibrary{decorations}
\usetikzlibrary{decorations.markings}
\usetikzlibrary{decorations.pathmorphing}
\usetikzlibrary{arrows.meta}


% allow figures to be placed in the running text by default:
\usepackage{scrhack}
\usepackage{float}
\floatplacement{figure}{htbp}
\floatplacement{table}{htbp}

% keep figures and tables in the section
\usepackage[section, below]{placeins}


%------------------------------------------------------------------------------
%---------------------- customize list environments ---------------------------
%------------------------------------------------------------------------------

\usepackage{enumitem}

%------------------------------------------------------------------------------
%------------------------------ Bibliographie ---------------------------------
%------------------------------------------------------------------------------

\usepackage[
  backend=biber,   % use modern biber backend
  autolang=hyphen, % load hyphenation rules for if language of bibentry is not
  sorting=none,    % german, has to be loaded with \setotherlanguages
                   % in the references.bib use langid={en} for english sources
]{biblatex}
\addbibresource{references.bib}  % die Bibliographie einbinden
\DefineBibliographyStrings{english}{andothers = {{et\,al\adddot}}} 

%------------------------------------------------------------------------------
%------------------------------ Sonstiges: ------------------------------------
%------------------------------------------------------------------------------

\usepackage[pdfusetitle,unicode,linkbordercolor=tugreen]{hyperref}
\usepackage{bookmark}
\usepackage[shortcuts]{extdash}
\usepackage{todo}

%------------------------------------------------------------------------------
%-------------------------    Angaben zur Arbeit   ----------------------------
%------------------------------------------------------------------------------

\author{Anja Beck}
\title{Flavour Mixing Effects in the Direct Detection of Dark Matter}
\date{2017}
\birthplace{Kempten (Allgäu)}
\chair{Lehrstuhl für Theoretische Physik IV}
\division{Fakultät Physik}
\thesisclass{Bachelor of Science}
\submissiondate{17. Juli 2017}
\firstcorrector{Jun.-Prof.~Dr.~Joachim Brod}
\secondcorrector{Prof.~Dr.~Heinrich Päs}

% tu logo on top of the titlepage
\titlehead{\includegraphics[height=1.5cm]{logos/tu-logo.pdf}}

\begin{document}
\frontmatter
%\input{content/hints.tex}
\maketitle

% Gutachterseite
\makecorrectorpage

% hier beginnt der Vorspann, nummeriert in römischen Zahlen
\thispagestyle{plain}
\section*{Abstract}
Quark flavour mixing is often neglected because of its small effects. But since dark matter only interacts weakly with baryonic matter, these tiny effects might be of importance for direct detection. This thesis examines whether flavour mixing was rightfully neglected by Altmannshofer et. al. in the discussion of a new interaction \cite{Z}. Therefore, we expand an existing framework for dark matter direct detection with the CKM mixing matrix. In the second part of this thesis, we present the new interaction, that was originally proposed to explain anomalies in the decay $B\rightarrow Kl\bar{l}$, but also makes predictions about the direct detection of dark matter. However, when treating dark matter cross sections, this new set-up ignores flavour mixing effects. Thus, we compare those new cross sections with cross sections obtained with flavour mixing and conclude that, in this case, flavour mixing was rightfully neglected.
\tableofcontents

\mainmatter
% Hier beginnt der Inhalt mit Seite 1 in arabischen Ziffern
\chapter{Introduction}
We will take a look at the effects of the flavour mixing mechanism on the direct detection of dark matter. Therefore we expand an existing formalism for dark matter direct detection with the CKM matrix. Afterwards we present a new interaction proposed to explain anomalies in the decay $B\rightarrow Kl\bar{l}$ that also makes predictions about the direct detection of dark matter. Finally, we compare these predictions with the direct detection cross sections that respect flavour mixing.

Direct detection means detecting dark matter directly through interaction with a nucleus, in contrary to indirect detection, which means measuring secondary products of dark matter annihilation or dark matter decay.

\todo{Irgendwo muss noch kurz auf dunkle Materie eingegangen werden. Als Einleitung wären daher auch Detection Experimente gut. Evtl. hier direct vs indirect erklären.}


\begin{figure}
	\centering
	\begin{tikzpicture}
\tikzstyle{centerArrow}=[decoration={
    markings,
    mark=at position 0.5 with {\fill (2pt,0)--(-2pt,2.31pt)--(-2pt,-2.31pt)--cycle;}}]
\begin{scope}
\def\xmove{2.5}
\def\ymove{1.25}
\def\centerSize{0.15}
\node [fill, circle,inner sep=\centerSize cm] (tCenter)  {};
\node (upperLeft) at (-\xmove,\ymove) {$\chi$};
\node (upperRight) at (\xmove,\ymove) {$\chi$};
\node (lowerLeft) at (-\xmove,-\ymove) {$N$};
\node (lowerRight) at (\xmove,-\ymove) {$N$};
\draw [centerArrow,postaction={decorate}]  (upperLeft) -- (tCenter) ;
\draw [centerArrow,postaction={decorate}]  (lowerLeft) -- (tCenter) ;
\draw [centerArrow,postaction={decorate}]  (tCenter) -- (upperRight) ;
\draw [centerArrow,postaction={decorate}]  (tCenter) -- (lowerRight) ;
\end{scope}
\end{tikzpicture}
	\caption{Direct detection}
	\label{fig:DirectDetection}
\end{figure}
\chapter{Introduction of the Flavour Mixing into an Existing Formalism}
\cite[Chapter 20]{Peskin}
\cite[Chapter 1.2.1]{Tevatron}
The origins of the flavour mixing go back to the 1960s, when Italian physicist Cabibbo resolved anomalies in weak decay data by proposing a flavour mixing of lefthanded down-type quarks. Later in 1973, Kobayaski and Maskawa extended this idea to three quark generation to explain CP violation. On a theoretical level the flavour mixing arises from the fact that the fermion mass eigenstates do not necessarily equal the flavour eigenstates. In the course of this chapter we introduce the Glashow-Weinberg-Salam Model of electroweak interaction and derive how the fermions gain their masses and why this leads to flavour mixing.

\section{The Higgs Mechanism}
We consider a complex scalar field $\Phi$ that interacts with itself through a potential
\begin{align}
	V(\Phi) = -\mu^2\Phi^\dagger\Phi + \frac{\lambda}{2}(\Phi^\dagger\Phi)^2 \ , \quad \mu^2 > 0 \ .
\end{align}
If $\lambda>0$, the minimum of this potential occurs at
\begin{align}
	\langle\Phi\rangle = \Phi_0 =  \sqrt{\frac{\mu^2}{\lambda}} \ .
\end{align}
This is the vacuum expectation value of $\Phi$.

\subsubsection{$SU(2)$}
We add a $SU(2)$ gauge field coupled to $\Phi$, so $\Phi$ is a doublet $(\Phi_1,\Phi_2)$ with covariant derivative
\begin{align}
	D_\mu\Phi = (\partial_\mu - ig\sum_{a=1}^3A_\mu^a\tau^a)\Phi \ ,
\end{align}
where $\tau^a$ are the generators of $SU(2)$. In this case there is an infinite number of vacuum expection values for $\Phi$ arranged in a circle. We are free to choose one, and we make the simple choice
\begin{align}
	\langle\Phi\rangle = \frac{1}{\sqrt{2}}\begin{pmatrix} 0 \\ v \end{pmatrix} \ .
\end{align}
Note that by choosing one vacuum value, we break the symmetry.


The kinetic energy of $\Phi$ is
\begin{align}
	(D_\mu\Phi)^2 = &\frac{1}{2}(\partial_\mu v)(\partial^\mu v)\notag \\
	- &ig(\partial^\mu\begin{pmatrix} 0 & v \end{pmatrix}) \left(\sum_{a=1}^3A_\mu^a\tau^a\begin{pmatrix} 0 \\ v \end{pmatrix}\right)\notag \\
	- &\frac{1}{2}g^2\begin{pmatrix} 0 & v \end{pmatrix}\sum_{a,b=1}^{3}\tau^a\tau^b\begin{pmatrix} 0 \\ v \end{pmatrix}A_\mu^aA^{b\mu} \ .
\end{align}
Using the relation $\{\tau^a,\tau^b\}=\sfrac{1}{2}\cdot\delta_{ab}$, we can simplify the last expression to get
\begin{align}
	- \frac{1}{2}g^2\begin{pmatrix} 0 & v \end{pmatrix}\sum_{a,b=1}^{3}\tau^a\tau^b\begin{pmatrix} 0 \\ v \end{pmatrix}A_\mu^aA^{b\mu} &= - \frac{g^2v^2}{8}\sum_{a=1}^{3}A_\mu^aA^{a\mu} \ ,
\end{align}
which is a mass term $\mathcal{L}_m = -\frac{1}{2}m_A^2A_\mu A^\mu$ that assigns the mass $m_A = \frac{gv}{2}$ to all three gauge bosons.

\subsubsection{$SU(2)\times U(1)$}
We expand the system with an additional $U(1)$ symmetry with gauge boson $B$. The field $\Phi$ has a charge $Y_\Phi$ under $U(1)$. The new covariant derivative is
\begin{align}
	D_\mu\Phi = \left(\partial_\mu - ig\sum_{a=1}^3A_\mu^a\tau^2 - iY_\Phi g'B_\mu\right)\Phi \ .
\end{align}
Again, we examine the kinetic term
\begin{align}
	(D_\mu\Phi)^2 = &\frac{1}{2}(\partial_\mu v)(\partial^\mu v)\notag \\
	- &\frac{1}{2}i\left(\partial^\mu\begin{pmatrix} 0 & v \end{pmatrix}\right) \left(g\sum_{a=1}^3A_\mu^a\tau^2 + Y_\Phi g'B_\mu\right)\begin{pmatrix} 0 \\ v \end{pmatrix}\notag \\
	- &\frac{1}{2}\begin{pmatrix} 0 & v \end{pmatrix}
	\left(g^2\sum_{a,b=1}^3A_\mu^aA^{b\mu}\tau^a\tau^b + 2gg'Y_\Phi\sum_{a=1}^3A_\mu^a\tau^aB^\mu + Y_\Phi^2g'^2B_\mu B^\mu\right)
	\begin{pmatrix} 0 \\ v \end{pmatrix}
\end{align}
Using $\{\tau^a,\tau^b\}=\sfrac{1}{2}\cdot\delta_{ab}$ and replacing $\tau^a=\sfrac{\sigma^a}{2}$, we find for the last term
\begin{align}
	\mathcal{L}^{(\text{mass})} &= - \frac{v^2}{2}\left(g^2\frac{1}{4}\sum_{a=1}^3A_\mu^aA^{a\mu} + Y_\Phi^2g'^2B_\mu B^\mu - gg'Y_\Phi B^\mu A_\mu^3\right)\notag \\
	&= - \frac{1}{2}\frac{v^2}{4}\left(g^2A_\mu^1A^{1\mu} + g^2A_\mu^2A^{2\mu} + (gA_\mu^3 - 2g'Y_\Phi B_\mu)^2\right)\notag \\
	&= - \frac{1}{2}\frac{v^2}{4}\left(g^22W^+W^- + (g^2+4g'^2Y_\Phi^2)Z_0^2\right) \ .
\end{align}
Here we identified the known vector bosons
\begin{align}
	W_\mu^\pm &= \frac{1}{\sqrt{2}}(A_\mu^1\mp iA_\mu^2) \ , &&\quad &m_W &= \frac{v}{2}g \notag \\
	Z_\mu^0 &= \frac{1}{\sqrt{g^2+4g'^2Y_\Phi^2}}(gA_\mu^3 - 2g'Y_\Phi B_\mu) \ , &&\quad &m_Z &= \frac{v}{2}\sqrt{g^2+4g'^2Y_\Phi^2}
\end{align}
as mass eigenstates of the gauge bosons.


Since we are in a $SU(2)\times U(1)$ symmetry, there has to be a fourth gauge boson. As we have just derived, it is massless. We renounce giving an elaborate explanation for this, but we want to give a motivation. Therefore we need to remember that the masses arise through choosing a vacuum expectation value for $\Phi$ and thereby breaking the $SU(2)$ symmetry. Looking at the gauge transformation of $\Phi$
\begin{align}
	\Phi\rightarrow e^{i\sum_{a=1}^3\alpha^a\tau^a}e^{i\beta Y_\Phi}\Phi \ ,
\end{align}
we find that the choice $\alpha^1=\alpha^2=0$, $\alpha^3 = 2\beta Y_\Phi$ leaves the vacuum expectation value unchanged. Thus, parts of the symmetry are conserved and keep one gauge boson from acquiring mass. The fourth gauge boson is the photon, and it is orthogonal to $Z_\mu^0$:
\begin{align}
	A_\mu = \frac{1}{\sqrt{g^2 + 2g'^2Y_\Phi^2}}(2g'Y_\Phi A^3_\mu + gB_\mu) \ .
\end{align}

\section{Fermion Masses and Flavour Mixing}
In this and the next chapters we use the following notation for the chiral components of the particle functions. For the leptons:
\begin{align}
	E_R &= (e_R,\mu_R,\tau_R) \ , &&\quad &Y_E &= -1 \ ; \\
	L_L &= \left(\begin{pmatrix} \nu_e \\ e \end{pmatrix}_L,
	\begin{pmatrix} \nu_\mu \\ \mu \end{pmatrix}_L,
	\begin{pmatrix} \nu_\tau \\ \tau \end{pmatrix}_L\right) \ , &&\quad &Y_L &= -\frac{1}{2} \ ;
\end{align}
and accordingly for the quarks
\begin{align}
	U_R &= (u_R,c_R,t_R) \ , &&\quad &Y_U &= \frac{2}{3} \ ; \\
	D_R &= (d_R,s_R,b_R) \ , &&\quad &Y_D &= -\frac{1}{3} \ ; \\
	Q_L &= \left(\begin{pmatrix} u \\ d \end{pmatrix}_L,
	\begin{pmatrix} c \\ s \end{pmatrix}_L,
	\begin{pmatrix} t \\ b \end{pmatrix}_L\right) \ , &&\quad &Y_Q &= \frac{1}{6} \ .	
\end{align}
The charge under $U(1)_Y$, also hypercharge, is $Y$. It is related to the electric charge $Q$ and the third component of the weak isospin $I_3$ through the Gell-Mann–Nishijima formula: $Y = 2(Q-I_3)$. The righthanded particles are singlets under $SU(2)$ and therefore $I_3=0$. We will use as well $E_L = (e_L,\mu_L,\tau_L)$, $D_L = (d_L,s_L,b_L)$, and $U_L = (u_L,c_L,t_L)$.



The electroweak interaction lagrangian for the standard model fermions is
\begin{align}\label{eq:LFermion}
%	\mathcal{L}^{(\text{int})} = &\sum_{i=1}^{3}\bar{E}_R^i(i\slashed{\partial}-g_1Y_E\slashed{B})E_R^i + \bar{D}_R^i(i\slashed{\partial}-g_1Y_D\slashed{B})D_R^i + \bar{U}_R^i(i\slashed{\partial}-g_1Y_U\slashed{B})U_R^i \notag \\
%	+ &\sum_{i=1}^{3}\bar{L}_L^i(i\slashed{\partial}-g_1Y_L\slashed{B}-g_2\slashed{A})L_L^i + \bar{Q}_L^i(i\slashed{\partial}-g_1Y_Q\slashed{B}-g_2\slashed{A})Q_L^i \ ,
	\mathcal{L}^{(\text{int})} = &\bar{E}_R(i\slashed{\partial}-g_1Y_E\slashed{B})E_R + \bar{D}_R(i\slashed{\partial}-g_1Y_D\slashed{B})D_R + \bar{U}_R(i\slashed{\partial}-g_1Y_U\slashed{B})U_R \notag \\
	+ &\bar{L}_L(i\slashed{\partial}-g_1Y_L\slashed{B}-g_2\slashed{A})L_L + \bar{Q}_L(i\slashed{\partial}-g_1Y_Q\slashed{B}-g_2\slashed{A})Q_L \ .
\end{align}



The lagrangian in \eqref{eq:LFermion} describes massless particles. A fermion mass term couples the lefthanded and righthanded part of a particle. Since direct coupling between for example $e_R$ and $(\nu_e,e)_L$ would violate gauge invariance, a connecting field. To preserve invariance under Lorentz, $U(1)_Y$, and $SU(2)$ transformations this field must have spin 0, hypercharge $Y=\sfrac{1}{2}$, and be a doublet. We identify this field with $\Phi$ from the previous chapter and write the mass terms for the fermions
\begin{align}
%	\mathcal{L}^{(\text{mass})} = -\sum_{i,j=1}^{3}\left[\lambda^e_{ij} \bar{L}_L^i\Phi E_R^j  + \lambda_{ij}^d\bar{Q}_L^i\Phi D_R^j + \lambda_{ij}^u\sum_{a,b=1}^{2}\epsilon_{ab}\bar{Q}_{La}^{i}\Phi_b^\dagger U_R^j + \text{h.c.}\right] \ ,
	\mathcal{L}^{(\text{mass})} = -\left[\bar{L}_L\Phi\lambda^e  E_R  + \bar{Q}_L\Phi\lambda^d D_R + \bar{Q}_Li\sigma^2\Phi^\dagger\lambda^u U_R + \text{h.c.}\right] \ ,
\end{align}
with complex matrix coupling constants $\lambda^e,\lambda^d,\lambda^u$. Replacing $\Phi$ with its vacuum expectation value gives
\begin{align}\label{eq:LMass}
%	\mathcal{L}^{(\text{mass})} = -\frac{v}{\sqrt{2}}\sum_{i,j=1}^{3}\left[\lambda^e_{ij} \bar{E}_L^i E_R^j  + \lambda_{ij}^d\bar{D}_L^i D_R^j + \lambda_{ij}^u\bar{U}_L^{i} U_R^j + \text{h.c.}\right] \ ,
	\mathcal{L}^{(\text{mass})} = -\frac{v}{\sqrt{2}}\left[ \bar{E}_L\lambda^e E_R  + \bar{D}_L \lambda^dD_R + \bar{U}_L\lambda^u U_R + \text{h.c.}\right] \ .
\end{align}



The interaction lagrangian \eqref{eq:LFermion} is invariant under unitary transformations
\begin{align}
	E_L &\rightarrow S_eE_L && &E_R &\rightarrow R_eE_R \ , \\
	U_L &\rightarrow S_uU_L && &U_R &\rightarrow R_uU_R \ , \\
	D_L &\rightarrow S_dD_L && &D_R &\rightarrow R_dD_R \ .
\end{align}
Thus, we can diagonalize the interactions in \eqref{eq:LMass}. The diagonal lepton coupling is $\tilde{\lambda}^e = S_e\lambda^eR_e^\dagger$ and parametrizes the lepton masses
\begin{align}
	m_e = \frac{v}{\sqrt{2}}\tilde{\lambda}_{11}^e \ , \quad m_\mu = \frac{v}{\sqrt{2}}\tilde{\lambda}_{22}^e \ , \quad m_\tau = \frac{v}{\sqrt{2}}\tilde{\lambda}_{33}^e \ .
\end{align}
The diagonal coupling for up-type quarks is $\tilde{\lambda}^u = S_u\lambda^uR_u^\dagger$, giving the corresponding masses
\begin{align}
	m_u = \frac{v}{\sqrt{2}}\tilde{\lambda}_{11}^u \ , \quad m_c = \frac{v}{\sqrt{2}}\tilde{\lambda}_{22}^u \ , \quad m_t = \frac{v}{\sqrt{2}}\tilde{\lambda}_{33}^u \ .
\end{align}
The transformed coupling of the down-type quarks is $\tilde{\lambda}^d = S_d\lambda^dR_d^\dagger$, leading to the down-type masses
\begin{align}
	m_d = \frac{v}{\sqrt{2}}\tilde{\lambda}_{11}^d \ , \quad m_s = \frac{v}{\sqrt{2}}\tilde{\lambda}_{22}^d \ , \quad m_b = \frac{v}{\sqrt{2}}\tilde{\lambda}_{33}^d \ .
\end{align}


The transformed particle functions are now mass eigenstates. But when looking at couplings of up- and down-type quarks, e.g. the current
\begin{align}
	\bar{U}_L\gamma^\mu D_L \ ,
\end{align}
the transformation changes the interaction to
\begin{align}
	\bar{U}_L\gamma^\mu S_u^\dagger S_d D_L \ .
\end{align}
The matrix $V = S_u^\dagger S_d$ is the CKM matrix.
In this chapter we include the flavour mixing into an existing formalism. We use the model described in \cite{ChiralEFT}. It provides a framework to calculate cross sections for the direct detection of dark matter. The model bases on a set of dimension-five, -six, and -seven operators. We restrict our calculations to the dimension-six operators, which are
\begin{align}\label{eq:normal}
	R_{1,q} &= (\bar{\chi}\gamma_\mu\chi)(\bar{q}\gamma^\mu q) && &R_{3,q} &= (\bar{\chi}\gamma_\mu\chi)(\bar{q}\gamma^\mu\gamma_5q)\notag \\
	R_{2,q} &= (\bar{\chi}\gamma_\mu\gamma_5\chi)(\bar{q}\gamma^\mu q) &&	&R_{4,q} &= (\bar{\chi}\gamma_\mu\gamma_5\chi)(\bar{q}\gamma^\mu\gamma_5q) \ .
\end{align}
Since the CKM mixing only applies to the lefthanded down-type quarks, we need to rewrite these operators in terms of the left- and righthanded particle functions to include the CKM matrix. These chiral operators are
\begin{align}\label{eq:chiral}
	Q_{1ij} &= (\bar{\chi}\gamma_\mu\tilde{\tau}^a\chi)(\bar{Q}_L^i\gamma^\mu \tau^aQ_L^j) && &Q_{5ij} &= (\bar{\chi}\gamma_\mu\gamma_5\tilde{\tau}^a\chi)(\bar{Q}_L^i\gamma^\mu \tau^aQ_L^j)\notag \\
	Q_{2ij} &= (\bar{\chi}\gamma_\mu\chi)(\bar{Q}_L^i\gamma^\mu Q_L^j) && &Q_{6ij} &= (\bar{\chi}\gamma_\mu\gamma_5\chi)(\bar{Q}_L^i\gamma^\mu Q_L^j)\notag \\
	Q_{3ij} &= (\bar{\chi}\gamma_\mu\chi)(\bar{u}_R^i\gamma^\mu u_R^j) && &Q_{7ij} &= (\bar{\chi}\gamma_\mu\gamma_5\chi)(\bar{u}_R^i\gamma^\mu u_R^j)\notag \\
	Q_{4ij} &= (\bar{\chi}\gamma_\mu\chi)(\bar{d}_R^i\gamma^\mu d_R^j) && &Q_{8ij} &= (\bar{\chi}\gamma_\mu\gamma_5\chi)(\bar{d}_R^i\gamma^\mu d_R^j) \ ,
\end{align}
here $Q_L^i = (u_L^i, d_L^i)$ is the isospin doublet of the $i^\text{th}$ quark generation. The operators $\tilde{\tau}^a,\tau^a$ are the generators of the SU(2) in the corresponding spin-representation. So for the quarks, we use
\begin{align}
	\tau^a = \frac{\sigma_a}{2} \ ,
\end{align}
where $\sigma_a$ are the pauli matrices. Regarding the dark matter particle, we decide to only keep the terms with $\tilde{\tau}^a = \tilde{\tau}^3$ and get
\begin{align}
	\tilde{\tau}^3\chi = \tau_0\chi \ ,
\end{align}
with the weak isospin value $\tau_0$ of the dark matter particle.

Including the CKM matrix in the fromalism follows these steps:
\begin{enumerate}
	\item Replace the pure lefthanded down-type quarks with the mixed quarks.
	\item Rewrite the chiral particle functions in terms of the normal particle functions and projection operators.
	\item Write down the entire interaction lagrangian in terms of \eqref{eq:normal} and \eqref{eq:chiral} and compare the coefficients.
\end{enumerate}

\section{Including the Flavour Mixing}
The inclusion of the CKM matrix only affects the chiral operators $Q_{1ij},Q_{2ij},Q_{5ij},Q_{6ij}$. The quark part of the interaction with flavour mixing is therefore
\begin{align*}
	\bar{Q}_L^i\gamma^\mu Q_L^j &= \begin{pmatrix}
	\bar{u}_L^i \\ \bar{d}_L^i
	\end{pmatrix}
	\gamma^\mu \begin{pmatrix}
	u_L^j \\ d_L^j
	\end{pmatrix}
	= \bar{u}_L^i\gamma^\mu d_L^j + \bar{d}_L^i\gamma^\mu d_L^j \\
	&= \bar{u}_L^i\gamma^\mu u_L^j + (V_{id}^*\bar{d}_L + V_{is}^*\bar{s}_L+V_{ib}^*\bar{b}_L)\gamma^\mu(V_{jd}d_L+V_{js}s_L+V_{jb}b_L)
\end{align*}
for $Q_{2ij}, Q_{6ij}$, respectively
\begin{align*}
	\bar{Q}_L^i\gamma^\mu\tau^a Q_L^j &=  \begin{pmatrix}
	\bar{u}_L^i \\ \bar{d}_L^i
	\end{pmatrix}
	\gamma^\mu \begin{pmatrix}
	u_L^j \\ d_L^j
	\end{pmatrix}
	= \frac{1}{2}\bar{u}_L^i\gamma^\mu \sigma_a d_L^j + \frac{1}{2}\bar{d}_L^i\gamma^\mu\sigma_a d_L^j \\
	&= \frac{1}{2}\bar{u}_L^i\gamma^\mu \sigma_a u_L^j + \frac{1}{2}(V_{id}^*\bar{d}_L + V_{is}^*\bar{s}_L+V_{ib}^*\bar{b}_L)\gamma^\mu\sigma_a(V_{jd}d_L+V_{js}s_L+V_{jb}b_L)
\end{align*}
for $Q_{1ij}, Q_{5ij}$. For simplicity we only keep the light quarks $u,d,s$ and neglect mixed terms. So the whole set of quark interactions is
\begin{align}
	\bar{Q}_L^i\gamma^\mu\tau^a Q_L^j &\approx \frac{1}{2}\bar{u}_L\gamma^\mu u_L\delta_{ij}\delta_{iu}\delta_{a3} - \frac{1}{2}\delta_{a3}(V_{id}^*V_{jd}\bar{d}_L\gamma^\mu d_L + V_{is}^*V_{js}\bar{s}_L\gamma^\mu s_L)\notag \\
	\bar{Q}_L^i\gamma^\mu Q_L^j &\approx \bar{u}_L\gamma^\mu u_L\delta_{ij}\delta_{iu} + V_{id}^*V_{jd}\bar{d}_L\gamma^\mu d_L + V_{is}^*V_{js}\bar{s}_L\gamma^\mu s_L\notag \\
	\bar{u}_R^i\gamma^\mu u_R^j &\approx\bar{u}_R\gamma^\mu u_R\delta_{ij}\delta_{iu}\notag \\
	\bar{d}_R^i\gamma^\mu d_R^j &\approx \bar{d}_R\gamma^\mu d_R\delta_{ij}\delta_{id} + \bar{s}_R\gamma^\mu s_R\delta_{ij}\delta_{is} \ .
\end{align}

\section{Replacing Chiral Particle Functions}
The next step is rewriting the chiral particle functions in terms of the normal particle functions using the projection operators $P_L,P_R$. This leads us to
\begin{align}
	\bar{Q}_L^i\gamma^\mu\tau^a Q_L^j =&\frac{1}{4}(\bar{u}\gamma^\mu u\delta_{ij}\delta_{iu}\delta_{3a}
	- V_{id}^*V_{jd}\bar{d}\gamma^\mu d\delta_{3a} - V_{is}^*V_{js}\bar{s}\gamma^\mu s\delta_{3a})\notag \\
	-&\frac{1}{4} (\bar{u}\gamma^\mu \gamma_5u\delta_{ij}\delta_{iu}\delta_{3a} - V_{id}^*V_{jd}\bar{d}\gamma^\mu\gamma_5 d\delta_{3a} - V_{is}^*V_{js}\bar{s}\gamma^\mu\gamma_5 s\delta_{3a})\notag \\
	%%%%%%%%%%%%%%%%%%%%%%%%%%%%%%%%%%%%%%%%%%%%%%%%%%%%%%%%%%%%%%%%%%%%%%%%%%%%%%%%%%%%%%%%%%%%%%%%%%%%%%%%%%%%%%%%%%%%
	\bar{Q}_L^i\gamma^\mu Q_L^j =&\frac{1}{2}(
	\bar{u}\gamma^\mu u\delta_{iu}\delta_{ij} + V_{id}^*V_{jd}\bar{d}\gamma^\mu d
	+ V_{is}^*V_{js}\bar{s}\gamma^\mu s)\notag \\
	-&\frac{1}{2}(\bar{u}\gamma^\mu\gamma_5u\delta_{iu}\delta_{ij} + V_{id}^*V_{jd}\bar{d}\gamma^\mu\gamma_5d + V_{is}^*V_{js}\bar{s}\gamma^\mu\gamma_5s)\notag \\
	%%%%%%%%%%%%%%%%%%%%%%%%%%%%%%%%%%%%%%%%%%%%%%%%%%%%%%%%%%%%%%%%%%%%%%%%%%%%%%%%%%%%%%%%%%%%%%%%%%%%%%%%%%%%%%%%%%%%
	\bar{u}_R^i\gamma^\mu u_R^j =&\frac{1}{2}(\bar{u}\gamma^\mu u\delta_{ij}\delta_{iu} + \bar{u}\gamma^\mu \gamma_5u\delta_{ij}\delta_{iu})\notag \\
	%%%%%%%%%%%%%%%%%%%%%%%%%%%%%%%%%%%%%%%%%%%%%%%%%%%%%%%%%%%%%%%%%%%%%%%%%%%%%%%%%%%%%%%%%%%%%%%%%%%%%%%%%%%%%%%%%%%%%
	\bar{d}_R^i\gamma^\mu d_R^j =&\frac{1}{2}(\bar{d}\gamma^\mu d\delta_{ij}\delta_{id} + \bar{d}\gamma^\mu\gamma_5d\delta_{ij}\delta_{id} + \bar{s}\gamma^\mu s\delta_{ij}\delta_{is} + \bar{s}\gamma^\mu \gamma_5s\delta_{ij}\delta_{is}) \ .
\end{align}
At this point we can express the chiral operators in terms of the original operators from \eqref{eq:normal}:
\begin{align}\label{eq:dependenciesoperators}
	Q_{1ij} =&\frac{\delta_{3a}\tau_{0}}{4}(R_{1u}\delta_{ij}\delta_{iu} - V_{id}^*V_{jd}R_{1d} - V_{is}^*V_{js}R_{1s})\notag \\
	-&\frac{\delta_{3a}\tau_{0}}{4} (R_{3u}\delta_{ij}\delta_{iu} - V_{id}^*V_{jd}R_{3d} - V_{is}^*V_{js}R_{3s})\notag \\
	%%%%%%%%%%%%%%%%%%%%%%%%%%%%%%%%%%%%%%%%%%%%%%%%%%%%%%%%%%%%%%%%%%%%%%%%%%%%%%%%%%%%%%%%%%%%%%%%%%%%%%%%%%%%%%%%%%%%%%%%
	Q_{2ij} =&\frac{1}{2}(R_{1u}\delta_{iu}\delta_{ij}+ V_{id}^*V_{jd}R_{1d}+ V_{is}^*V_{js}R_{1s})\notag \\
	-&\frac{1}{2}(R_{3u}\delta_{iu}\delta_{ij} + V_{id}^*V_{jd}R_{3d} + V_{is}^*V_{js}R_{3s})\notag \\
	%%%%%%%%%%%%%%%%%%%%%%%%%%%%%%%%%%%%%%%%%%%%%%%%%%%%%%%%%%%%%%%%%%%%%%%%%%%%%%%%%%%%%%%%%%%%%%%%%%%%%%%%%%%%%%%%%%%%%%%
	Q_{3ij} =&\frac{1}{2}(R_{1u}\delta_{ij}\delta_{iu} + R_{3u}\delta_{ij}\delta_{iu})\notag \\
	%%%%%%%%%%%%%%%%%%%%%%%%%%%%%%%%%%%%%%%%%%%%%%%%%%%%%%%%%%%%%%%%%%%%%%%%%%%%%%%%%%%%%%%%%%%%%%%%%%%%%%%%%%%%%%%%%%%%%%%
	Q_{4ij} =&\frac{1}{2}(R_{1d}\delta_{ij}\delta_{id} + R_{3d}\delta_{ij}\delta_{id} + R_{1s}\delta_{ij}\delta_{is} + R_{3s}\delta_{ij}\delta_{is}) \ .
	%%%%%%%%%%%%%%%%%%%%%%%%%%%%%%%%%%%%%%%%%%%%%%%%%%%%%%%%%%%%%%%%%%%%%%%%%%%%%%%%%%%%%%%%%%%%%%%%%%%%%%%%%%%%%%%%%%%%%%%
%	Q_{5ij} &=\frac{\delta_{3a}\tau_{0}}{4}(R_{2u}\delta_{ij}\delta_{iu}
%	- V_{id}^*V_{jd}R_{2d} - V_{is}^*V_{js}R_{2s})\notag \\
%	&-\frac{\delta_{3a}\tau_{0}}{4} (R_{4u}\delta_{ij}\delta_{iu}
%	- V_{id}^*V_{jd}R_{4d} - V_{is}^*V_{js}R_{4s})\notag \\
%	%%%%%%%%%%%%%%%%%%%%%%%%%%%%%%%%%%%%%%%%%%%%%%%%%%%%%%%%%%%%%%%%%%%%%%%%%%%%%%%%%%%%%%%%%%%%%%%%%%%%%%%%%%%%%%%%%%%%%%%%
%	Q_{6ij} &= \frac{1}{2}(R_{2u}\delta_{iu}\delta_{ij} + V_{id}^*V_{jd}R_{2d} + V_{is}^*V_{js}R_{2s})\notag \\
%	&- \frac{1}{2}(R_{4u}\delta_{iu}\delta_{ij} + V_{id}^*V_{jd}R_{4d} + V_{is}^*V_{js}R_{4s})\notag \\
%	%%%%%%%%%%%%%%%%%%%%%%%%%%%%%%%%%%%%%%%%%%%%%%%%%%%%%%%%%%%%%%%%%%%%%%%%%%%%%%%%%%%%%%%%%%%%%%%%%%%%%%%%%%%%%%%%%%%%%%%
%	Q_{7ij} &= \frac{1}{2}(R_{2u}\delta_{ij}\delta_{iu} + R_{4u}\delta_{ij}\delta_{iu})\notag \\
%	%%%%%%%%%%%%%%%%%%%%%%%%%%%%%%%%%%%%%%%%%%%%%%%%%%%%%%%%%%%%%%%%%%%%%%%%%%%%%%%%%%%%%%%%%%%%%%%%%%%%%%%%%%%%%%%%%%%%%%%
%	Q_{8ij} &= \frac{1}{2}(R_{2d}\delta_{ij}\delta_{id} + R_{4d}\delta_{ij}\delta_{id} + R_{2s}\delta_{ij}\delta_{is} + R_{4s}\delta_{ij}\delta_{is}) \ .
\end{align}
The operators $Q_{5ij}-Q_{8ij}$ can be obtained from $Q_{1ij}-Q_{4ij}$ by replacing $R_{1q}\leftrightarrow R_{2q}$ and $R_{3q}\leftrightarrow R_{4q}$.

\section{Comparing Coefficients}
Our final step is expressing the coefficients $K_{l,q}$ of the original operators $R_{l,q}$ in \eqref{eq:normal} in terms of the coefficients $C_{lij}$ of the chiral operators $Q_{lij}$ in \eqref{eq:chiral}. To get there we look at the overall interaction. The interaction cannot depend on the representation of particle functions we choose, either normal or chiral, therefore
\begin{align}
	\sum_{l,q} K_{l,q}R_{l,q} \overset{!}{=} \sum_{l,i,j}C_{lij}Q_{lij} \ .
\end{align}
By putting the interactions \eqref{eq:dependenciesoperators} into the right side of the equation and rearranging the expression in terms of the $R_{l,q}$, we conclude that $K_{l,q}$ must be equal to the terms in front of $R_{l,q}$ on the right side. We get the dependencies
\begin{align}\label{eq:CoeffNormal}
	K_{1,u} &= \sum_{i,j}\frac{\delta_{ij}\delta_{iu}}{2}\left(C_{1ij}\frac{\delta_{3a}\tau_{0}}{2}+ 
	C_{2ij} + C_{3ij}\right)\notag \\
	K_{1,d} &= \sum_{i,j}\frac{1}{2}\left(-V_{id}^*V_{jd}C_{1ij}\frac{\delta_{3a}\tau_{0}}{2}+ V_{id}^*V_{jd}C_{2ij} + \delta_{ij}\delta_{id}C_{4ij}\right)\notag \\
	K_{1,s} &= \sum_{i,j}\frac{1}{2}\left(- V_{is}^*V_{js}C_{1ij}\frac{\delta_{3a}\tau_{0}}{2}+ V_{is}^*V_{js}C_{2ij}+ \delta_{ij}\delta_{is}C_{4ij}\right)\notag \\
	%%%%%%%%%%%%%%%%%%%%%%%%%%%%%%%%%%
	K_{2,u} &= \sum_{i,j}\frac{\delta_{ij}\delta_{iu}}{2}\left(\frac{\delta_{3a}\tau_{0}}{2}C_{5ij}+ 
	C_{6ij} + C_{7ij}\right)\notag \\
	K_{2,d} &= \sum_{i,j}\frac{1}{2}\left(- V_{id}^*V_{jd}\frac{\delta_{3a}\tau_{0}}{2}C_{5ij} + C_{6ij}V_{id}^*V_{jd} + \delta_{ij}\delta_{id}C_{8ij}\right)\notag \\
	K_{2,s} &= \sum_{i,j}\frac{1}{2}\left(- V_{is}^*V_{js}\frac{\delta_{3a}\tau_{0}}{2}C_{5ij} + C_{6ij}V_{is}^*V_{js} + \delta_{ij}\delta_{is}C_{8ij}\right)\notag \\
	%%%%%%%%%%%%%%%%%%%%%%%%%%%%%%%%%%
	K_{3,u} &= \sum_{i,j}\frac{\delta_{ij}\delta_{iu}}{2}\left(- C_{1ij}\frac{\delta_{3a}\tau_{0}}{2} - C_{2ij} + C_{3ij}\right)\notag \\
	K_{3,d} &= \sum_{i,j}\frac{1}{2}\left(C_{1ij}\frac{\delta_{3a}\tau_{0}}{2} V_{id}^*V_{jd} - V_{id}^*V_{jd}C_{2ij} + \delta_{ij}\delta_{id}C_{4ij}\right)\notag \\
	K_{3,s} &= \sum_{i,j}\frac{1}{2}\left(C_{1ij}\frac{\delta_{3a}\tau_{0}}{2} V_{is}^*V_{js} - V_{is}^*V_{js}C_{2ij} + \delta_{ij}\delta_{is}C_{4ij}\right)\notag \\
	%%%%%%%%%%%%%%%%%%%%%%%%%%%%%%%%%%
	K_{4,u} &= \sum_{i,j}\frac{\delta_{ij}\delta_{iu}}{2}\left(-\frac{\delta_{3a}\tau_{0}}{2}C_{5ij} - C_{6ij} + C_{7ij}\right)\notag \\
	K_{4,d} &= \sum_{i,j}\frac{1}{2}\left(\frac{\delta_{3a}\tau_{0}}{2}C_{5ij}V_{id}^*V_{jd} - C_{6ij}V_{id}^*V_{jd} + \delta_{ij}\delta_{id}C_{8ij}\right)\notag \\
	K_{4,s} &= \sum_{i,j}\frac{1}{2}\left(\frac{\delta_{3a}\tau_{0}}{2}C_{5ij}V_{is}^*V_{js} - C_{6ij}V_{is}^*V_{js} + \delta_{ij}\delta_{is}C_{8ij}\right) \ .
\end{align}
To have a hermitian interaction the coefficients need to fulfil the relation $C_{lij} = C_{lji}^*$.
\chapter{The \texorpdfstring{$L_\mu-L_\tau$}{TEXT} Model}
In this chapter we present an extension to the standard model proposed in \cite{InColour}. The authors originally aimed at explaining anomalies in the decay $B\rightarrow Kl\bar{l}$, but also obtained predictions for the direct detection of dark matter in the succeeding publication \cite{Z}. We will later compare their results with the formalism in the previous chapter including the CKM mixing.

\section{The New Interaction}
The extension to the standard model in \cite{InColour} is a new $U(1)'$ gauge group. The related vector-boson is called $Z'$, and it couples to the muon and tau lepton families, and a new set of vector-like quarks $U,D,Q$. The standard model quarks indirectly couple to the $Z'$ as well, since they mix with the new quarks through a Yukawa coupling:
\begin{align}
	\mathcal{L}^{(\text{mix})} \ .
\end{align}
In \cite{Z} also a coupling to a dark matter fermion $\chi$ is established. Before present the entire interaction, we want to mention that the $Z'$ and the new quarks get their masses through a new higgs-like field $\Phi$. Its vacuum expectation value is $\langle\Phi\rangle=\sfrac{v_\Phi}{\sqrt{2}}$, and it connects the coupling strength $g'$ to the $Z'$ mass: $m_{Z'} = v_\Phi g'$. The full interaction lagrangian is:
\begin{align}
	\mathcal{L}^{(\text{int})}_{Z'} = &g'Z_\alpha'\times q_l\left(\bar{l}_2\gamma^\alpha l_2 - \bar{l}_3\gamma^\alpha l_3 + \bar{\mu}_R\gamma^\alpha\mu_R-\bar{\tau}_R\gamma^\alpha\tau_R\right)\notag \\
	+&g'Z_\alpha'\times v_\Phi^2\left(-\frac{Y_{Di}Y^*_{Dj}}{2m_D^2}\bar{d}_R^i\gamma^\alpha d_R^j - \frac{Y_{Ui}Y^*_{Uj}}{2m_U^2}\bar{u}_R^i\gamma^\alpha u_R^j + \bar{u}_L^i\gamma^\alpha u_L^j + \bar{d}_L^i\gamma^\alpha d_L^j\right)\notag \\
	+&g'Z_\alpha'\times q_\chi(\bar{\chi}\gamma^\alpha\chi) \ ,
\end{align}
where $q_l,q_\chi$ are the $U(1)'$ charge of the leptons and the dark matter particle, $l_2,l_3$ are the electroweak lefthanded lepton doublets, $u^i,d^i$ are the standard model quarks, and $m_{U,D,Q}$ are the masses of the new quarks.


\todo{Was genau sind die neuen Quarks? Gibt es davon nur 3?}
\todo{Erkläre die Mischung der neuen und SM Quarks etwas ausführlicher.}
\todo{Kann man tatsächlich hier die Ups drehen und dann später mit gedrehten Downs rechnen? Die entsprechenden Konstanten dann noch einfügen.}

\section{Restrictions to the Parameter Space}
In \cite{Z} Altmannshofer et. al. discuss restrictions for the parameter space by looking at the $B$ decay mentioned above and discussing dark matter relic density and direct detection. They find that experimental data from $B\rightarrow Kl\bar{l}$ limits the ratio of the $Z'$ mass and the coupling $g'$ to
\begin{align}
	\SI{540}{\giga\electronvolt}\lessapprox\frac{m_{Z'}}{g'}\lessapprox\SI{4.9}{\tera\electronvolt} \ ,
\end{align}
with $m_{Z'}\gtrapprox\SI{10}{\giga\electronvolt}$.


Regarding the dark matter relic density, they conclude that only
\begin{align}
	m_{Z'}\approx 2m_\chi
\end{align}
leads to correct results. Since they neglect flavour mixing in the nucleus, direct detection has to occur through the loop diagram in figure \ref{fig:Loop}. The corresponding cross section at zero momentum transfer is (see \cite{Z})
\begin{align}
	\sigma_{0,\text{Loop}} = \frac{\mu_{A\chi}^2}{A^2\pi}\left(\frac{\alpha_{em}Z}{3\pi}\ \frac{g'^2q_\chi q_l}{m_{Z'}^2}\log\left(\frac{m_\mu^2}{m_\tau^2}\right)\right)^2 \ ,
\end{align}
where $\mu_{A\chi}$ is the reduced mass of the nucleus and the dark matter particle $\chi$ and $A,Z$ are the nucleon and proton numbers.


When discussing limits to the parameter space they distinguish two cases. For $q_l=q_\chi=1$, experimental data favours the parameter region
\begin{align}
	\SI{10}{\giga\electronvolt}\lessapprox m_\chi\lessapprox\SI{46}{\giga\electronvolt}\notag \\
	\SI{2e-3}{}\lessapprox g'\lessapprox\SI{e-2}{} \ ,
\end{align}
leaving possible dark matter masses in the range $(5-23)\si{\giga\electronvolt}$. For $q_l=1,q_\chi=\sfrac{1}{6}$ no further restriction of the parameters can be found.
\begin{figure}
	\centering
	\begin{tikzpicture}
\tikzstyle{centerArrow}=[decoration={
	markings,
	mark=at position 0.5 with {\fill (2pt,0)--(-2pt,2.31pt)--(-2pt,-2.31pt)--cycle;}}]

\begin{scope}[xshift=6cm,yshift=5cm]
\def\xmove{2}
\def\ymove{1.25}
\def\centerShift{2cm}
\def\centerCircle{1cm}
\def\centerSize{0.05cm}
\coordinate (tCenter1) at (0,0);
\coordinate (tCenter2) at (0,-\centerShift);
\coordinate (tCenter3) at (0,-\centerShift-\centerCircle);
\coordinate (tCenter4) at (0,-\centerShift-\centerCircle-\centerShift);

\node (upperLeft) at (-\xmove,\ymove) {$\chi$};
\node (upperRight) at (\xmove,\ymove) {$\chi$};
\node (lowerLeft) at (-\xmove,-\centerShift-\centerCircle-\centerShift-\ymove cm) {$N$};
\node (lowerRight) at (\xmove,-\centerShift-\centerCircle-\centerShift-\ymove cm) {$N$};

\draw [centerArrow,postaction={decorate}]  (upperLeft) -- (tCenter1) ;
\draw [centerArrow,postaction={decorate}]  (tCenter1) -- (upperRight) ;
\draw [centerArrow,postaction={decorate}]  (lowerLeft) -- (tCenter4) ;
\draw [centerArrow,postaction={decorate}]  (tCenter4) -- (lowerRight) ;
\draw [decoration={snake, segment length=1.5mm, amplitude=0.5mm},decorate] (tCenter1) -- (tCenter2) ;
\node at (0.5,-\centerShift/2) {$Z'$};
\draw [
        decoration={markings, mark=at position 0.5 with {\fill (2pt,0)--(-2pt,2.31pt)--(-2pt,-2.31pt)--cycle;}, mark=at position 1 with {\fill (2pt,0)--(-2pt,2.31pt)--(-2pt,-2.31pt)--cycle;}},
        postaction={decorate}
] ([yshift=.5cm]tCenter3) ellipse(.45 and 0.5);
\node at ([yshift=-\centerCircle/2,xshift=0.75cm]tCenter2) {$l$};
\node at ([yshift=-\centerCircle/2,xshift=-0.75cm]tCenter2) {$l$};
\draw [decoration={snake, segment length=1.5mm, amplitude=0.5mm},decorate] (tCenter3) -- (tCenter4) ;
\node at (0.5,-\centerShift/2-\centerCircle-\centerShift) {$\gamma$};
\end{scope}
\end{tikzpicture}
	\caption{Direct detection loop diagram}
	\label{fig:Loop}
\end{figure}

\todo{Evtl. kann man den Ursprung der Grenzen noch näher erläutern.}
\chapter{Comparison}
\section{Cross Section}
We are only interested in the interactions $Q_{2bs}$, therefore all coefficients $C_{lij}$ vanish except $C_{2sb} = C_{2bs}^*$. In terms of the coefficients in \eqref{eq:CoeffNormal} this means
\begin{align}
	K_{1,d} &= +\text{Re}(V_{cd}^*V_{td}C_{2sb})\notag \\
	K_{1,s} &= +\text{Re}(V_{cs}^*V_{ts}C_{2sb})\notag \\
	K_{3,d} &= -\text{Re}(V_{cd}^*V_{td}C_{2sb})\notag \\
	K_{3,s} &= -\text{Re}(V_{cs}^*V_{ts}C_{2sb}) \ ,
\end{align}
and all the other $K_{l,q}$ are zero. \\
We are interested in spin-independent cross sections, therefore we only use the interactions $(\bar{\chi}\gamma^\mu\chi)(\bar{q}\gamma_\mu q)$. Protons do not contain strange quarks, so the lagrangian only consists of one interaction:
\begin{align}
	\mathcal{L} = K_{1,d}(\bar{\chi}\gamma^\mu\chi)(\bar{d}\gamma_\mu d) + \text{h.c.} \ .
\end{align}
Since this operator only counts the number of down quarks in the nucleus, the matrix element is
\begin{align}
	M = Z\cdot 2K_{1,d} + (A-Z)\cdot K_{1,d} \ .
\end{align}
\begin{align}
	\sigma_{0,\text{tree}}^\text{SI} &= \frac{\mu_{A\chi}^2}{A^2\pi}\left|ZC_p +(A-Z)C_n\right|^2 \ . \\
	&= \frac{\mu_{A\chi}^2}{A^2\pi}K_{1,d}^2\times\mathcal{O}(10^2)
\end{align}



In case of the spin-dependent cross section we have
\begin{align}
	\sigma_{0,\text{tree}}^\text{SD} &= \frac{\mu_{A\chi}^2}{A^2\pi}32\Lambda^2J(J+1)64G_F^2 \\
	&= \frac{\mu_{A\chi}^2}{A^2\pi}32\left(\frac{K_{3,d}}{J}\frac{\Delta d^{(n)}}{\sqrt{2}G_F}\frac{\mu}{3.826}\right)^2J(J+1)64G_F^2 \\
	&= \frac{\mu_{A\chi}^2}{A^2\pi}2^{10}K_{3,d}^2(\Delta d^{(n)})^2\left(\frac{\mu}{3.826}\right)^2\frac{J(J+1)}{J^2} \\
	&= \frac{\mu_{A\chi}^2}{A^2\pi}K_{3,d}^2\times\mathcal{O}(10^{3}10^{-2}10^{-2})
\end{align}

\begin{align}
	&\langle N|K_{1,d}\bar{d}\gamma^\mu d|N\rangle + \langle N|K_{1,u}\bar{u}\gamma^\mu u|N\rangle\notag \\
	= &Z\left(2K_{1,u}+K_{1,d}\right)(\bar{p}\gamma^\mu p) + (A-Z)\left(K_{1,u}+2K_{1,d}\right)(\bar{n}\gamma^\mu n)\notag \\
	\notag \\
	&\langle N|K_{3,d}\bar{d}\gamma^\mu\gamma_5 d|N\rangle + \langle N|K_{3,u}\bar{u}\gamma^\mu\gamma_5 u|N\rangle\notag \\
	= &Z(K_{3,d}2s^\mu\Delta d^{(p)} + K_{3,u}2s^\mu\Delta u^{(p)})(\bar{p}\gamma^\mu\gamma_5 p)\notag \\
	+ &(A-Z)(K_{3,d}2s^\mu\Delta d^{(n)} + K_{3,u}2s^\mu\Delta u^{(n)})(\bar{n}\gamma^\mu\gamma_5 n)\notag
\end{align}
\begin{figure}
	\centering
	\begin{tikzpicture}
\tikzstyle{centerArrow}=[decoration={
markings,
mark=at position 0.5 with {\fill (2pt,0)--(-2pt,2.31pt)--(-2pt,-2.31pt)--cycle;}}]
\begin{scope}[xshift=3cm,yshift=-2cm]
\def\xmove{2}
\def\ymove{1.25}
\def\centerShift{2}
\def\centerSize{0.08cm}
\coordinate (tCenter1) at (0,0);
\coordinate[fill, circle,inner sep=\centerSize] (tCenter2) at (0,-\centerShift cm);
\node (upperLeft) at (-\xmove,\ymove) {$\chi$};
\node (upperRight) at (\xmove,\ymove) {$\chi$};
\node (lowerLeft) at (-\xmove,-\centerShift cm-\ymove cm) {$b_L,s_L$};
\node (lowerRight) at (\xmove,-\centerShift cm-\ymove cm) {$s_L,b_L$};
\node at (0.5,-\centerShift/2) {$Z'$};
\draw [centerArrow,postaction={decorate}]  (upperLeft) -- (tCenter1) ;
\draw [centerArrow,postaction={decorate}]  (tCenter1) -- (upperRight) ;
\draw [centerArrow,postaction={decorate}]  (lowerLeft) -- (tCenter2) ;
\draw [centerArrow,postaction={decorate}]  (tCenter2) -- (lowerRight) ;
\draw [decoration={snake, segment length=1.5mm, amplitude=0.5mm},decorate] (tCenter1) -- (tCenter2) ;
\end{scope}
\end{tikzpicture}
	\caption{Direct detection tree level diagram}
	\label{fig:DD}
\end{figure}
\todo{Wieso ist eigentlich $Q_L^2\gamma_\mu Q_L^3 = s_L\gamma_\mu b_L$?}
\chapter{Conclusion}
\input{content/conclusion.tex}

\appendix
% Hier beginnt der Anhang, nummeriert in lateinischen Buchstaben


\todos
\backmatter
\printbibliography

\cleardoublepage
\thispagestyle{empty}
\section*{Eidesstattliche Versicherung}
Ich versichere hiermit an Eides statt, dass ich die vorliegende Abschlussarbeit mit dem Titel \enquote{\thetitle} selbstständig und ohne unzulässige fremde Hilfe erbracht habe.
Ich habe keine anderen als die angegebenen Quellen und Hilfsmittel benutzt, sowie wörtliche und sinngemäße Zitate kenntlich gemacht. 
Die Arbeit hat in gleicher oder ähnlicher Form noch keiner Prüfungsbehörde vorgelegen.

\vspace*{1cm}\noindent
\begin{center}
  \begin{tabular}{@{}p{0.4\textwidth}@{\hspace{0.15\textwidth}}p{0.4\textwidth}@{}}
  \rule{\linewidth}{0.25pt}& \rule{\linewidth}{0.25pt}\\
  Ort, Datum & Unterschrift
  \end{tabular}
\end{center}

\subsection*{Belehrung}
Wer vorsätzlich gegen eine die Täuschung über Prüfungsleistungen betreffende Regelung einer Hochschulprüfungsordnung verstößt, handelt ordnungswidrig.
Die Ordnungswidrigkeit kann mit einer Geldbuße von bis zu \SI[round-mode=places, round-precision=2]{50000}{€} geahndet werden. 
Zuständige Verwaltungsbehörde für die Verfolgung und Ahndung von Ordnungswidrigkeiten ist der Kanzler/die Kanzlerin der Technischen Universität Dortmund. 
Im Falle eines mehrfachen oder sonstigen schwerwiegenden Täuschungsversuches kann der Prüfling zudem exmatrikuliert werden (\S\,63 Abs. 5 Hochschulgesetz --HG--).

Die Abgabe einer falschen Versicherung an Eides statt wird mit Freiheitsstrafe bis zu 3 Jahren oder mit Geldstrafe bestraft.

Die Technische Universität Dortmund wird ggf.\ elektronische Vergleichswerkzeuge (wie z.\,B.\ die Software \enquote{turnitin}) zur Überprüfung von Ordnungswidrigkeiten in Prüfungsverfahren nutzen. \\[\baselineskip]

\noindent Die oben stehende Belehrung habe ich zur Kenntnis genommen.\\[1cm]
\begin{center}
\begin{tabular}{@{}p{0.4\textwidth}@{\hspace{0.15\textwidth}}p{0.4\textwidth}@{}}
\rule{\linewidth}{0.25pt}& \rule{\linewidth}{0.25pt}\\
Ort, Datum & Unterschrift
\end{tabular}
\end{center}

\end{document}
